\section{Задание}

\subsection{Цель работы}
	Изучить возможности системы управления версиями и коллективной
	разработкой программного проекта; освоить методы работы с утилитой Git
	для управления версиями и коллективной разработкой программного проекта.

\subsection{Средства выполнения}
\begin{itemize}
	\item утилита Git;
	\item утилита Git GUI и gitk.
\end{itemize}

\subsection{Задачи}

В ходе выполнения лабораторной работы были выполнены следующие пункты задания:
\begin{itemize}
	\item запустить Git GUI и создать новый репозиторий в рабочей папке проекта;
	\item добавить в папку репозитория файл, зафиксировать состояние репозитория (выполнить commit);
	\item внести изменения в файл, зафиксировать новое состояние репозитория;
	\item создать новую ветку (lastname), внести в неё изменения и зафиксировать их;
	\item переключиться на ветку master, внести в неё изменения (исправление) и зафиксировать их;
	\item продемонстрировать слияние веток - разрешить возникший конфликт;
	\item просмотреть дерево изменений веток (историю), в том числе на примере выноса части кода в отдельные функцию и файлы;
	\item продемонстрировать откат изменений (revert);
	\item создать удалённый репозиторий на github.com и опубликовать в нём изменения.

\end{itemize}
